\section{花絮、阴谋论与被低估的玩家}

这一节收录了访谈中零散但极有价值的洞察和故事——它们不属于任何单一主题，但每一个都揭示了 AI 产业的某个侧面。

\subsection{阴谋论 vs 真相}

围绕 DeepSeek 的阴谋论在社交媒体上疯传。Dylan 和 Nathan 逐一拆解：

\textbf{``中国政府在补贴 DeepSeek''}：不太可能。DeepSeek 由幻方量化（对冲基金）资助，不是政府关联企业。梁文锋后来确实与政府领导会面，但那是 DeepSeek 成功\textbf{之后}——政府是在蹭热度，不是在背后资助。

\textbf{``他们在发布前做空了 NVIDIA 股票''}：也不太可能。V3 在 12 月 26 日发布——谁会选择圣诞节后第一天发布来配合做空？\textit{``我觉得他们只是在赶工——谁在乎圣诞节，赶在中国新年前发出来就行。''} (Nathan)

\textbf{``只花了 600 万美元训练''}：这是对论文的严重误读。\$5.576M 是 V3 预训练\textbf{一次运行}的 GPU 小时成本。不包括前期实验（可能跑了数十次小规模消融）、后训练、R1 训练、推理服务。实际总投入可能是这个数字的 10-50 倍。

\subsection{被低估的 Gemini Flash Thinking}

Lex 的个人测试中，Google 的 Gemini Flash Thinking 可能比 R1 更便宜且不弱——但几乎没人谈论它。原因可能是：
\begin{itemize}
    \item Google 的营销远不如 DeepSeek 的开源策略吸引眼球
    \item Flash Thinking 可能使用了不同的方法（在现有架构上叠加推理，而非专门的推理训练）
    \item DeepSeek 的``中国黑马''叙事更有传播力
\end{itemize}

这是一个关于``技术 vs 叙事''的教训——最好的技术不一定赢得最多关注。

\subsection{PyTorch PowerPlant.NoBlowUp}

训练大型集群时，GPU 的功率波动可以摧毁电力设备。计算阶段高功率（所有 GPU 全力计算）$\to$ 通信阶段低功率（GPU 等待数据交换）。这种脉冲式负载对变压器和配电设备非常危险。

Meta 的工程师写了一个 PyTorch 算子来解决这个问题：在 GPU 空闲期让它们计算无意义的数字（假乘法），纯粹为了维持稳定功率。这段代码意外被开源——一个叫 \texttt{PowerPlant.no\_blowup = 1} 的操作符。

\textit{``你很容易就能把东西炸了。''} (Dylan) 这个故事完美说明了 AI 训练中``无聊但关键''的工程挑战——问题不是算法突破，而是如何不让你的电力变压器爆炸。

\subsection{Anthropic 的安全困境}

Anthropic 据报道拥有比 O3 更好的推理模型但因安全考虑不发布。R1 的 Chain-of-Thought 确实可以令人不安——在中英文之间切换、出现类似乱码的片段、然后突然给出正确答案。

DeepSeek 的激进发布降低了所有人的安全标准——类似于冷战时期苏联太空计划对美国 NASA 的压力。\textit{``DeepSeek ships. That's one of their big advantages.''} (Nathan) 快速发布本身就是一种竞争武器。

\subsection{Sam Altman 的预言}

引用 Sam Altman 的一句被频繁提到的话：\textit{``超人说服力将在超人智能之前到来。''} 这对 AI 安全的含义深远——在 AI 能``真正思考''之前，它可能已经能``真正说服''。Character AI 的聊天机器人已经在影响年轻用户的情绪——如果是故意的文化/政治操纵呢？

Zuckerberg 在财报电话会上的公开声明也耐人寻味：\textit{``对我们的国家优势来说，开源标准应该是美国的，这很重要。''} 开源不仅是技术策略，也是地缘政治工具。

\subsection{本章小结}

AI 产业充满了引人入胜的细节：从``不要炸了变压器''的工程挑战，到``600 万美元''的误读，到被低估的 Gemini，到开源作为地缘政治武器。这些花絮揭示了一个比技术论文更丰富、更混乱、更人性的行业现实。

